\appendix

\begin{center}
\Large
\textbf{MindForge: Empowering Embodied Agents with Theory of Mind for Lifelong
Cultural Learning} \\
\vspace{0.2em}Appendices \\
\vspace{0.3em}
\end{center}

\section{Additional Experimental Results}
\label{appendix:additional-experiments}

\subsection{Collaborative Experiments}
\label{appendix:comm-experiments}

We provide an extended set of experiments to reinforce our findings in \Cref{fig:comm-rounds} and \Cref{fig:weak-weak}. Specifically, we consider atomic tasks that are significantly harder than dirt and wood collection in order to demonstrate the generality of our framework: (1) mining an iron block and (2) crafting a wooden pickaxe.

\begin{table}[htb]
    \scriptsize
  \caption{Agent performance increases as the number of rounds grows in both collaborative and instructive settings. In the instructive setting, we use Llama 3.1-70B as the weak agent and GPT-4 as the strong agent. Similarly, collaborative experiments use Llama 3.1-70B for both \collabvoyager agents.}
    {\tabcolsep=9pt\def\arraystretch{1.1}
    \begin{tabularx}{\columnwidth}{@{}>{\arraybackslash}l *5{X}}
      \toprule
      \textbf{Setting} & \textbf{Task} & \textbf{Comm. Round 0} &\textbf{Comm. Round 1} & \textbf{Comm. Round 2} & \textbf{Comm. Round 3} \\ 
      \midrule
      Instructive & Mine Iron & $41\%$ & $50\%$ & $54\%$ & $62\%$ \\
      Instructive & Craft a Pickaxe & $20\%$ & $33\%$ & $41\%$ &$45\%$ \\
      \midrule
      Collaborative & Craft a Pickaxe & $20\%$ & $25\%$ & $33\%$ & $33\%$ \\
      \bottomrule
    \end{tabularx}}
    \label{table:comm-experiments-appendix}
\end{table}

As \Cref{table:comm-experiments-appendix} shows, extend communication and perspective-taking between \collabvoyager agents enable significant improvements on complex Minecraft tasks that require compositionality ($41\% \rightarrow 62\%$ and $20\% \rightarrow 45\%$). Similarly to \Cref{fig:weak-weak}, we see that a purely collaborative setting results in smaller performance gains ($20\% \rightarrow 33\%$) than instructive learning, due to reinforcing wrong beliefs and similar thinking patterns.

\subsection{Post-Collaboration and OOD tasks}
\label{appendix:post-collab}

Do social interactions have long-term effects on agents and help them perform better on different tasks? We consider the impact of interactions in post-collaboration and out-of-distribution scenarios to assess \collabvoyager's generalization and lifelong learning capabilities. Specifically, we define post-collaboration scenarios as a non-collaborative setting where the \collabvoyager agent tries to complete a task based on prior beliefs formed from a previous collaborative experience stored in the semantic memory. Moreover, we consider scenarios where the agent is tasked with solving a novel task with prior collaborative experience on an adjacent task in the Minecraft tech-tree. As \Cref{table:ood} shows, collaboration between \collabvoyager agents together with semantic memory showcase an average improvement of $8.33\%$ for in-distribution tasks (dirt collection) together with an average improvement of $8.78\%$ in non-collaborative out-of-distribution scenarios (wood collection). These results showcase the importance of semantic memory in achieving continual learning through collaboration. 

\begin{table}[h]
\caption{Post-collaboration and out-of-distribution generalization of the \collabvoyager agent. We report the task completion rates across $24$ individual trials.}
  \label{table:ood}
  \scriptsize
  {\tabcolsep=10pt\def\arraystretch{1.1}
  \begin{tabularx}{\columnwidth}{@{}>{\arraybackslash}l *4{X}}
  \toprule
  \textbf{Scenario} & \textbf{Task} & \textbf{Mistral-7B} & \textbf{Mixtral-8x7B} & \textbf{Llama 3.1-8B}\\ 
  \midrule
  No Collaboration & Dirt & $37.5\%$ & $29.15\%$ & $12.5\%$ \\
  Post-Collaboration & Dirt & \bm{$37.5\%$} & \bm{$41.66\%$} & \bm{$25\%$}\\
  \midrule
  No Collaboration & Wood & $33.3\%$ & $50\%$ & $7\%$ \\
  OOD post-collaboration & Wood & \bm{$41.66\%$} & \bm{$58.33\%$}& \bm{$16.66\%$}\\
  \bottomrule
  \end{tabularx}}
\end{table}

\subsection{Alternative Communication Protocols}
\label{appendix:additional}
\Cref{tab:flexible-communication} shows the performance of \collabvoyager agents when they are allowed to begin a conversation at point compared to when they can only initiate a conversation upon failure.

\begin{table}[hbt]
    \centering
    \caption{Comparison between the default \collabvoyager communication protocol and a more flexible alternative where agents can initiate the conversation at any point. Both protocols achieve similar results across different model sizes.}
    \label{tab:flexible-communication}
    \scriptsize
    \def\arraystretch{1.2}
    \begin{tabular}{@{}lllcccc@{}}
    \toprule
    \multirow{2}{*}{Model} & \multirow{2}{*}{Task} & \multirow{2}{*}{Setting} & \multicolumn{4}{c}{Task Completion Rate} \\
    \cmidrule(l){4-7}
    & & & Round 0 & Round 1 & Round 2 & Round 3 \\
    \midrule
    \multirow{4}{*}{Mixtral-8x7B} & \multirow{2}{*}{Mine dirt} & MindForge w/ flexible communication & 37\% & 45\% & 67\% & 67\% \\
    & & MindForge & 29\% & 42\% & 62\% & 67\% \\
    \cmidrule(l){2-7}
    & \multirow{2}{*}{Mine dirt and wood} & MindForge w/ flexible communication & 75\% & 79\% & 79\% & 83\% \\
    & & MindForge & 75\% & 79\% & 79\% & 83\% \\
    \midrule
    \multirow{4}{*}{Mistral-7B} & \multirow{2}{*}{Mine dirt} & MindForge w/ flexible communication & 37\% & 42\% & 45\% & 54\% \\
    & & MindForge & 37\% & 42\% & 45\% & 54\% \\
    \cmidrule(l){2-7}
    & \multirow{2}{*}{Mine dirt and wood} & MindForge w/ flexible communication & 41\% & 45\% & 50\% & 50\% \\
    & & MindForge & 41\% & 45\% & 50\% & 50\% \\
    \bottomrule
    \end{tabular}
\end{table}

\section{Additional Ablations}\label{appx:additional-ablations}
\subsection{Perspective Taking} \label{appendix:perspective}



To quantitatively assess the effect of perspective-taking, we perform an ablation study where we attempt to solve a Minecraft task without perspective-taking during communication in an instructive learning setting. \Cref{tab:perf-perspective} showcases how the absence of perspective-taking leads to worse performance, irrespective of how much agents communicate. Moreover, as the agents use more communication rounds when trying to solve a task, taking perspective of the other's agent situation increases task-completion rate.

\begin{table}[htbp]
    \centering
    \caption{Perspective-taking ablation. We consider the Minecraft task of collecting a dirt block. Improvement is quantified as the fraction of agents that solve the task across 24 individual trials.}
    \label{tab:perf-perspective}
    \scriptsize
    \def\arraystretch{1.2}
    \begin{tabular}{@{}lcccc@{}}
    \toprule
    \multirow{2}{*}{Model Variant} & \multicolumn{4}{c}{Task Completion Rate} \\
    \cmidrule(l){2-5}
    & Round 0 & Round 1 & Round 2 & Round 3 \\
    \midrule
    \collabvoyager w/ \textbf{perspective-taking} & 29\% & 42\% & 61\% & 67\% \\
    \collabvoyager w/o \textbf{perspective-taking} & 29\% & 37\% & 50\% & 54\% \\
    \midrule
    Improvement & 0\% & +5\% & +11\% & +13\% \\
    \bottomrule
    \end{tabular}
\end{table}

We attribute the positive correlation between performance and usage of perspective-taking during communication to the ability of the \collabvoyager agent to provide more relevant and directed advice. This observation is also supported by the substantial increase in performance in the first two rounds as presented in \Cref{tab:perf-perspective}, where the teacher agent provides more insightful information as it gets to know the other's agent situation better. 

Additionally, we further ablate our structured representation of Theory of Mind by considering an unstructured perspective-taking 2-step prompt presented in Think Twice \cite{wilf2023think}. \Cref{tab:perf-perspective-unstructured} showcases that a structured ToM representation provides a meaningful advantage on complex Minecraft tasks.

\begin{table}[htbp]
    \centering
    \caption{Structured perspective-taking ablation with unstructured alternative. The unstructured ToM approach is inspired from Think Twice \cite{wilf2023think}. Using a structured belief-system results in higher performance on complex Minecraft tasks.}
    \label{tab:perf-perspective-unstructured}
    \scriptsize
    \def\arraystretch{1.2}
    \begin{tabular}{@{}llcccc@{}}
    \toprule
    \multirow{2}{*}{Task} & \multirow{2}{*}{Model Variant} & \multicolumn{4}{c}{Task Completion Rate} \\
    \cmidrule(l){3-6}
    & & Round 0 & Round 1 & Round 2 & Round 3 \\
    \midrule
    Craft a Pickaxe & \collabvoyager & 20\% & \textbf{33}\% & \textbf{41}\% & \textbf{45}\% \\
    Craft a Pickaxe & \collabvoyager w/o structured ToM & 20\% & 29\% & 33\% & 41\% \\
    \midrule
    Mine Iron & \collabvoyager & 41\% & 50\% & 54\% & \textbf{62}\% \\
    Mine Iron & \collabvoyager w/o structured ToM & 41\% & 50\% & 54\% & 58\% \\
    \bottomrule
    \end{tabular}
\end{table}

\subsection{Memory Components} \label{appendix:memory}

\collabvoyager employs three types of memory subsystems, as discussed in \Cref{ssec:memory-subsys}. These components ensure that the agent can reuse and recall relevant information about the task it is trying to complete from past experiences and interactions with other agents. The importance of procedural memory with respect to task completion is already studied in \citep{wang2023voyager}, concluding that reusing previous code snippets becomes increasingly important with the difficulty of the tasks. To understand the effect of episodic memory in the \collabvoyager framework, we perform an ablation on the episodic memory component within two Minecraft tasks. \Cref{table:episodic-ablation} showcases the drop in completion rate without the use of the episodic memory component. While we do not consider the tech-tree experiments here, we expect that considering past failure episodes leads to increasingly better results.

\begin{table}[H]
\caption{Ablation on the episodic memory component in \collabvoyager on a single task. We report the task completion rates across 24 individual trials.}
\label{table:episodic-ablation}
\scriptsize
\def\arraystretch{1.2}
\begin{tabularx}{\columnwidth}{@{}>{\arraybackslash}l *3{Y}}
\toprule
\textbf{Scenario} & \textbf{Task} & \textbf{Mixtral-8x7B} & \textbf{Llama 3.1-8B}\\
\midrule
\collabvoyager w/ \textbf{episodic memory} & Dirt & $\bm{29.15\%}$ & $12.5\%$ \\ 
\collabvoyager w/o \textbf{episodic memory} & Dirt & $25\%$ & \bm{$16.66\%$}\\ 
\midrule
\collabvoyager w/ \textbf{episodic memory} & Wood & $\bm{50\%}$ & $\bm{7\%}$ \\ 
\collabvoyager w/o \textbf{episodic memory} & Wood & $45.83\%$& \bm{$7\%$}\\
\bottomrule
\end{tabularx}
\end{table}

\section{Implementation details}\label{appx:impl-details}

\subsection{Perspective Taking}
Perspective-taking represents a central part of our framework since it allows \collabvoyager agents to take advantage of the internal causal structure of other agents in order to aid more efficient and grounded communication. In practice, perspective-taking corresponds to an LLM call where the prompt contains instructions and relevant information required to understand the internal state of other \collabvoyager agents during communication. \Cref{fig:flow} showcases how perspective-taking integrates with the overall communication mechanism. 

\begin{figure}[H]
\centering
\includegraphics[width=0.8\columnwidth, clip]{figures/overall-flow.pdf}
\caption{Flow of \collabvoyager communication components used to generate a natural language response. The agents' internal process consists of two stages: (i) internal representation update and (ii) perspective-taking. Leveraging both processes the agent can then generate a natural language response and send it through the Minecraft chat.}
\label{fig:flow}
\end{figure}

Once the agent receives a message in the Minecraft chat, a first LLM call is executed where the agent updates the stateful internal representation about the agent that sent the message. This update process is shown in \Cref{fig:beliefs-evolution-rounds}. Subsequently, the updated internal representation is used as context together with the rest of the conversation to take perspective of the situation of the other \collabvoyager agent. Similar to the first stage, perspective-taking is implemented as a LLM call, with the following prompt:

\vspace{0.25cm}
\begin{tcolorbox}[title=Prompt for perspective-taking, label=box:prompt-perspective]
You are a Minecraft agent named \{\{name\}\} and you are having a conversation with another agent named \{\{other\_name\}\}. \\

Based on the current conversation and your knowledge about the other agent, \{\{other\_name\}\}, take the other agent's perspective to assess and describe your current understanding, knowledge state, and likely needs from \{\{other\_name\}\}'s perspective. \\

Here is the current conversation between you and \{\{other\_name\}\}:\{conversation\} \\

Here is your mental model of \{\{other\_name\}\}: \{world\_model\} \\

Perspective Analysis: 
\end{tcolorbox}

Below you can find an example of perspective-taking from Minecraft gameplay:

\begin{tcolorbox}[title=Expert agent taking perspective of weak agent, label=box:pers-prompt]
Weak’s current knowledge state includes understanding the biome they are in, the time of day, and the fact that there are
dark oak logs nearby. Weak knows that a wooden axe is necessary for mining logs, and that crafting this tool will be a
required step before proceeding with the task. Weak’s likely needs from Strong include further assistance in navigating the environment during the night to find a tree
or waiting until the day.
\end{tcolorbox}

\subsection{Episodic Memory}

We implement episodic memory as a Retrieval-Augmented Generation (RAG) using LangChain \citep{langchain2024}. Specifically, we embed all the episodes where the agent failed to complete the task using \textbf{text-embedding-ada-002}. An episode consists of the context used to generate the action, the code itself and the corresponding critic message. Subsequently, when the agent generates a new action, the RAG is queried to retrieve the \textbf{k=5} most relevant failure episodes with respect to the task at hand. Lastly, to ensure we do not populate the context window with too many tokens, we generate a summary (using an LLM) of the retrieved episodes and feed this summary in the context used to generate the action. Below you can find the corresponding prompts for generating the episodic summary:

\begin{tcolorbox}[title=System message, label=box:sys-episodic]
You are a helpful assistant tasked with summarizing past experience episodes and pointing out the causes of failure. Create a concise summary.
\end{tcolorbox}

\begin{tcolorbox}[title=Prompt message, label=box:prompt-episodic]
Please summarize these episodes and why they failed:\{combined\_episodes\}
\end{tcolorbox}

\subsection{Procedural Memory}

The procedural memory has an identical design and functionality with the skill library introduced in Voyager \citep{wang2023voyager}. Once a task gets successfully completed, we store the corresponding piece of code such that we can later reference it in adjacent tasks. This becomes increasingly important in tasks that require multiple steps, like crafting a pickaxe: the agent needs to first create a crafting table, create wooden sticks and only then attempt to craft the pickaxe. When we run the \collabvoyager framework in an open-ended setting, the procedural memory allows the agent to continuously evolve and create an increasingly large collection of reusable skills.

\subsection{Semantic Memory}

As described in \Cref{ssec:memory-subsys}, semantic memory stores certain beliefs the agent has about the environment, coming from experience with the environment and the other \collabvoyager agents. Whenever \collabvoyager attempts a new task, it creates a certain belief about how the task should be solved. For example, the agent might believe it requires tools to mine a block of dirt. This belief is then added to the context of the LLM when the environment action is generated. As highlighted in \Cref{sec:intro}, one of the primary failure cases of Voyager are the false beliefs that come from this LLM generation. However, compared to Voyager \citep{wang2023voyager}, we leverage the communication with other \collabvoyager agents such that these beliefs get updated and refined. Below we showcase an example of false belief correction when trying to collect a block of wood:



\begin{tcolorbox}[title=Initial belief inside the semantic memory, label=box:response-semantic-1]
\{"How to mine 1 wood log in Minecraft?": "Answer: To mine 1 wood log in Minecraft, you need to use an axe."\}
\end{tcolorbox}

\begin{tcolorbox}[title=Updated belief after communication, label=box:response-semantic-1]
\{"How to mine 1 wood log in Minecraft?": "Answer: To mine 1 wood log in Minecraft, you need to punch a tree with your bare hands."\}
\end{tcolorbox}
In addition to task-related memory, the \collabvoyager agent maintains a stateful set of beliefs about the other agents in the environment that are constantly updated across episodes. An example of this is presented in \Cref{fig:beliefs-evolution-rounds}, where the agent updates its representation based on the contents of the conversation. Thus, keeping track of various information and facts about other \collabvoyager agents acts as a semantic memory retrieval when agents communicate.



\section{Sample Conversation}
\label{appendix:sample-conv}

Each communication round between two \collabvoyager agents comprises of six sending messages in a turn-based manner through the Minecraft build-in chat. Below you can find a sample conversation showcasing how the expert agent (Strong) helps the weak (Weak) agent. 

\begin{tcolorbox}[title=Conversation about completing the task: "Mine 1 wood log"]

\textbf{Weak}: Hey, can you help me with Mine 1 wood log? \\

\textbf{Strong}: Absolutely, I can help you with that. Given your current location in the dark forest biome and the presence of dark oak logs around you, you should be able to mine a log without any issues. Here's the code you can use: `mine('dark\_oak\_log')`. This command will make your agent mine a dark oak log. Just make sure you are close enough to the tree when executing this. \\

\textbf{Weak}:  Thank you for the helpful guidance, Strong! However, in my case, I am currently in a dark forest biome as well, but I am having trouble finding a tree during the night. Could you suggest a way to find a tree during the day so I can mine the wood log? I also noticed that I don't have a wooden axe in my inventory, so I may need to craft one before attempting to mine the log. I appreciate your assistance in this matter. \\

\textbf{Strong}: Of course, I understand your predicament. Since it's currently night time, visibility might be low which makes it challenging to locate a tree. You can wait until daybreak for better visibility. Use the `wait\_until\_day()` function to wait until it's daytime. For crafting a wooden axe, you would need 3 wood planks and 2 sticks. You can obtain wood planks from a wood log by using the `craft('wooden\_plank', wood\_log)` command and sticks from 2 wooden planks using the `craft('stick', wooden\_plank)` command. Once you have these, you can craft a wooden axe using the `craft('wooden\_axe', wooden\_plank, stick)` function. Here's the code:

\begin{verbatim}
```
wait_until_day()
craft(`wooden_plank', wood_log)
craft(`stick', wooden_plank)
craft(`wooden_axe', wooden_plank, stick)
```
\end{verbatim}


Once you have your axe and it's day time, you should be able to easily locate and mine a tree for the wood log. \\ 

\textbf{Weak}:  Thank you, Strong. I appreciate your help.


                            
\end{tcolorbox}

\section{Fine-tuning setup}
\label{appendix:fine}
\subsection{Execution Logs}
For the first fine-tuning variant, we use execution logs generated by a Voyager agent powered by GPT-4. Specifically, these logs contain input and output text information obtained from the components specific to the architecture, such as the actor or the critic. We find this to be a meaningful way to transfer knowledge as well as reasoning capabilities from GPT-4 to Mistral-7B. To perform the fine-tuning, we use LoRA \citep{hu2021lora} due to its low computational requirements. 

\subsection{Wiki Data and Documentation}

The second fine-tuning variation we consider builds upon the initial setup presented above and incorporates more structured information about Minecraft by including the Minecraft Wiki \citep{fan2022minedojo} and Minecraft API documentation. Similar to the first variant we use LoRA \citep{hu2021lora}.

\section{Prompts}
\label{appendix:prompts}

\begin{tcolorbox}[title=Prompt for updating partner beliefs, label=box:response-semantic-1]
You are a Minecraft agent. \\ 

You just had a conversation with another agent based on a task you are trying to solve. \\

Based on the contents of the conversation and the previous beliefs, you have to create a set of beliefs that represent your perception of the other agent.
\end{tcolorbox}

\begin{tcolorbox}[title=Prompt for creating interaction beliefs, label=box:response-semantic-1]
You are a Minecraft agent. \\

You just had a conversation with another agent based on a task you are trying to solve. \\

Based on the contents of the conversation and the previous beliefs, you have to create a set of beliefs that that can help you complete the task.
\end{tcolorbox}

\vspace{0.25cm}
\begin{tcolorbox}[title=Prompt for perspective-taking, label=box:prompt-perspective]
You are a Minecraft agent named \{\{name\}\} and you are having a conversation with another agent named \{\{other\_name\}\}. \\

Based on the current conversation and your knowledge about the other agent, \{\{other\_name\}\}, take the other agent's perspective to assess and describe your current understanding, knowledge state, and likely needs from \{\{other\_name\}\}'s perspective. \\

Here is the current conversation between you and \{\{other\_name\}\}:\{conversation\} \\

Here is your mental model of \{\{other\_name\}\}: \{world\_model\} \\

Perspective Analysis: 
\end{tcolorbox}

\section{LLMs} \label{appendix:llm}
This section details the technical specifications of the Large Language Models (LLMs) and related infrastructure used to implement the \collabvoyager framework and conduct the experiments presented in this paper.

All LLM calls, unless explicitly stated otherwise, were executed via API calls to external services. The exception to this is the fine-tuned model, which was ran and trained locally.

The specific models and their corresponding access methods are outlined below:

\begin{itemize}
    \item \textbf{Llama 3.1-8B-Instruct}: This model was accessed through the Lambda Inference API.\footnote{\url{https://docs.lambdalabs.com/public-cloud/lambda-inference-api/}} This API provided a reliable and efficient interface for interacting with the Llama 3.1-8B-Instruct model, enabling seamless integration within the \collabvoyager framework.
    \item \textbf{GPT-4}:  Accessed via the OpenAI API. The OpenAI API offered access to the advanced capabilities of the GPT-4 model, crucial for establishing performance baselines and for the expert agent in collaborative settings.
    \item \textbf{Mistral-7B Instruct v0.2, Mixtral-8x7B Instruct and Llama 3.1-70B}: These Mistral models were accessed through the Together AI API.\footnote{\url{https://api.together.ai/models}} The Together AI API provided a platform for utilizing these open-weight models, allowing for a comparative analysis within the Voyager and \collabvoyager frameworks.
\end{itemize}